\section{Sujet n\textsuperscript{o}\,30}
\noindent\begin{minipage}{0.5\linewidth}\footnotesize Office du Baccalauréat du Cameroun\end{minipage}%
\hfill\begin{minipage}{0.45\linewidth}\footnotesize\raggedleft Examen : \textbf{Baccalauréat} · Série : \textbf{C/E}\\
Épreuve : \textbf{Mathématiques} · Durée : \textbf{4\,h} · Coef. : \textbf{7}\end{minipage}
\par\smallskip{\color{onbo}\rule{\linewidth}{1pt}}

\par\smallskip\noindent{\footnotesize\itshape Sujet de synthèse — contexte : la centrale
hydroélectrique du barrage de \textbf{Lom Pangar} (Est-Cameroun). Calculatrice non programmable autorisée.}

\subsection*{Partie A — Évaluation des ressources \hfill (15 points)}

\noindent\textbf{Exercice 1.} \hfill\textit{(3 points)}\\
Sur cinq saisons, on a relevé la pluviométrie $x$ (en centaines de \si{\milli\metre}) sur le bassin
versant du barrage et la production électrique $y$ correspondante (en \si{\giga\watt\hour}) :
\begin{center}\small
\begin{tabular}{|c|*{5}{c|}}\hline
Pluviométrie $x_i$ & 12 & 14 & 16 & 18 & 20\\\hline
Production $y_i$ & 23 & 26 & 30 & 33 & 38\\\hline
\end{tabular}
\end{center}
\begin{enumerate}[label=\arabic*)]
\item Calculer les coordonnées $(\bar x\,;\,\bar y)$ du point moyen $G$, puis le coefficient de
corrélation linéaire $r$ de cette série. Interpréter. \hfill\textit{(1,5 pt)}
\item Déterminer une équation de la droite de régression de $y$ en $x$ par la méthode des moindres
carrés, sous la forme $y=ax+b$. \hfill\textit{(1 pt)}
\item Pour une saison où la pluviométrie atteindrait \SI{2400}{\milli\metre}, estimer la production
électrique du barrage. \hfill\textit{(0,5 pt)}
\end{enumerate}

\smallskip
\noindent\textbf{Exercice 2.} \hfill\textit{(3 points)}\\
Les turbines de la centrale sont entretenues selon deux cycles décalés. La planification conduit à
l'équation $(E):\ 11x-7y=1$, d'inconnue $(x\,;\,y)\in\Z^2$.
\begin{enumerate}[label=\arabic*)]
\item Vérifier que $(2\,;\,3)$ est une solution de $(E)$, puis résoudre $(E)$ dans $\Z^2$. \hfill\textit{(1,5 pt)}
\item Un compteur d'impulsions affiche, après $2025$ démarrages, la valeur $11^{2025}$. Déterminer
le reste de la division euclidienne de $11^{2025}$ par $7$. \hfill\textit{(1,5 pt)}
\end{enumerate}

\smallskip
\noindent\textbf{Exercice 3.} \hfill\textit{(4 points)}\\
Le rendement instantané (normalisé) d'une turbine en fonction du débit $x\ge0$ est modélisé par la
fonction $f$ définie sur $[0\,;\,+\infty[$ par $f(x)=x\,\mathrm{e}^{\,1-x}$, de courbe $(\mathcal{C})$.
On considère par ailleurs, pour tout entier $n\ge0$, l'intégrale
$\displaystyle I_n=\int_0^{1}x^{\,n}\,\mathrm{e}^{-x}\dd x$.
\begin{enumerate}[label=\arabic*)]
\item Étudier les variations de $f$ sur $[0\,;\,+\infty[$ et déterminer $\displaystyle\lim_{x\to+\infty}f(x)$ ;
en déduire le débit pour lequel le rendement est maximal et la valeur de ce maximum. \hfill\textit{(1,25 pt)}
\item À l'aide d'une intégration par parties, montrer que pour tout réel $\lambda>0$,
$\displaystyle\int_0^{\lambda}f(x)\dd x=\mathrm{e}-(\lambda+1)\,\mathrm{e}^{\,1-\lambda}$, puis calculer la limite de cette
intégrale quand $\lambda\to+\infty$. \hfill\textit{(1,25 pt)}
\item Montrer que la suite $(I_n)$ est décroissante. \hfill\textit{(0,75 pt)}
\item Établir que pour tout $n\ge0$, $0\le I_n\le\dfrac{1}{n+1}$, puis en déduire $\displaystyle\lim_{n\to+\infty}I_n$. \hfill\textit{(0,75 pt)}
\end{enumerate}

\smallskip
\noindent\textbf{Exercice 4.} \hfill\textit{(5 points)}\\
Le plan complexe est muni d'un repère orthonormé direct $(O\,;\,\vec u,\vec v)$. On note
$P(z)=z^2-2\sqrt3\,z+4$.
\begin{enumerate}[label=\arabic*)]
\item Résoudre dans $\C$ l'équation $P(z)=0$. On note $z_A$ la solution de partie imaginaire positive
et $z_B$ l'autre. \hfill\textit{(1 pt)}
\item Écrire $z_A$ et $z_B$ sous forme trigonométrique, puis calculer $z_A^{\,6}$. \hfill\textit{(1 pt)}
\item Soient $A$ et $B$ les points d'affixes $z_A$ et $z_B$. Démontrer que le triangle $OAB$ est
équilatéral. \hfill\textit{(1 pt)}
\item Soit $S$ la similitude directe de centre $O$, de rapport $2$ et d'angle $\dfrac\pi3$. Donner
l'écriture complexe de $S$, puis déterminer l'affixe de l'image de $A$ par $S$. \hfill\textit{(1 pt)}
\item Dans l'espace rapporté à un repère orthonormé $(O\,;\,\vec\imath,\vec\jmath,\vec k)$, le mur du
barrage passe par les points $M(2\,;\,0\,;\,0)$, $N(0\,;\,3\,;\,0)$ et $Q(0\,;\,0\,;\,6)$. Déterminer
une équation cartésienne du plan $(MNQ)$, puis calculer la distance du point $O$ à ce plan. \hfill\textit{(1 pt)}
\end{enumerate}

\subsection*{Partie B — Évaluation des compétences \hfill (5 points)}

\noindent\textit{La société de gestion du barrage de Lom Pangar sollicite ton expertise d'élève de
Terminale~C pour trois décisions techniques. \textbf{(a)} L'évacuateur de crue a une section droite
rectangulaire ; pour limiter les frottements, on impose un périmètre mouillé (les deux parois
verticales de hauteur $h$ et le fond de largeur $b$) égal à \SI{12}{\metre}, soit $2h+b=12$. On veut
que la section laisse passer le débit maximal, ce qui revient à maximiser l'aire de la section.
\textbf{(b)} Lors de la première mise en eau, le volume stocké est de \SI{3,6}{\giga\metre\cubed} et
augmente de \SI{12}{\percent} chaque année ; la cote d'exploitation est atteinte à
\SI{6}{\giga\metre\cubed}. \textbf{(c)} Chacune des $5$ turbines tombe en panne, indépendamment des
autres, avec une probabilité de \SI{10}{\percent} sur une année ; on s'intéresse au nombre de
turbines défaillantes dans l'année.}

\begin{enumerate}[label=\textbf{Tâche \arabic*.},leftmargin=2.4em]
\item Déterminer les dimensions $h$ et $b$ de la section droite d'\emph{aire maximale}, ainsi que
cette aire. \hfill\textit{(1,5 pt)}
\item Déterminer l'année à partir de laquelle la cote d'exploitation de \SI{6}{\giga\metre\cubed}
sera atteinte. \hfill\textit{(1,5 pt)}
\item Déterminer la probabilité qu'\emph{au moins une} des $5$ turbines tombe en panne dans l'année,
puis le nombre moyen de turbines défaillantes attendu. \hfill\textit{(1,5 pt)}
\end{enumerate}
\par\smallskip{\footnotesize\itshape Présentation générale : 0,5 point.}

% ════════════════════════════════════════════════════
\corrige{du sujet 30}

\noindent\textbf{Partie A — Exercice 1.}
1) $\bar x=\dfrac{12+14+16+18+20}{5}=16$ et $\bar y=\dfrac{23+26+30+33+38}{5}=30$, donc $G(16\,;\,30)$.
On calcule les variances et la covariance : $V(x)=\dfrac{1}{5}\sum(x_i-\bar x)^2=\dfrac{40}{5}=8$,
$V(y)=\dfrac{138}{5}=27{,}6$ et $\mathrm{cov}(x,y)=\dfrac{1}{5}\sum(x_i-\bar x)(y_i-\bar y)=\dfrac{74}{5}=14{,}8$.
D'où $r=\dfrac{\mathrm{cov}(x,y)}{\sqrt{V(x)\,V(y)}}=\dfrac{14{,}8}{\sqrt{8\times27{,}6}}\approx0{,}996$.
Comme $r$ est très proche de $1$, la corrélation linéaire est excellente : un ajustement affine est justifié.

\noindent 2) $a=\dfrac{\mathrm{cov}(x,y)}{V(x)}=\dfrac{14{,}8}{8}=1{,}85$ et $b=\bar y-a\bar x=30-1{,}85\times16=0{,}4$.
La droite de régression de $y$ en $x$ a pour équation $\boxed{y=1{,}85\,x+0{,}4}$.

\noindent 3) $x=\SI{2400}{\milli\metre}$ correspond à $x=24$ (centaines de mm) :
$y=1{,}85\times24+0{,}4=44{,}8$. Production estimée : environ $\mathbf{44{,}8\ \text{GWh}}$.

\medskip
\noindent\textbf{Exercice 2.}
1) $11\times2-7\times3=22-21=1$ : $(2\,;\,3)$ est bien solution. Par différence avec $(E)$ :
$11(x-2)=7(y-3)$. Comme $11\wedge7=1$, le théorème de Gauss donne $7\mid x-2$, soit $x=2+7k$,
$k\in\Z$ ; en reportant, $11\cdot7k=7(y-3)$ d'où $y=3+11k$. Donc
$\mathcal{S}=\{(2+7k\,;\,3+11k),\ k\in\Z\}$.

\noindent 2) Modulo $7$ : $11\equiv4$. Les puissances de $4$ modulo $7$ sont $4^1\equiv4$,
$4^2\equiv16\equiv2$, $4^3\equiv8\equiv1\ [7]$, période $3$. Or $2025=3\times675$, donc
$11^{2025}\equiv4^{2025}\equiv(4^3)^{675}\equiv1\ [7]$. Le reste est $\mathbf{1}$.

\medskip
\noindent\textbf{Exercice 3.}
1) $f'(x)=\mathrm{e}^{\,1-x}+x\cdot(-\mathrm{e}^{\,1-x})=(1-x)\,\mathrm{e}^{\,1-x}$. Comme $\mathrm{e}^{\,1-x}>0$,
$f'(x)$ a le signe de $1-x$ : $f$ croît sur $[0\,;\,1]$, décroît sur $[1\,;\,+\infty[$. Le rendement
est donc maximal pour le débit $x=1$, et vaut $f(1)=1\cdot\mathrm{e}^{0}=1$. Par croissances comparées,
$\lim\limits_{x\to+\infty}x\,\mathrm{e}^{-x}=0$ donc $\lim\limits_{x\to+\infty}f(x)=\mathrm{e}\cdot0=0$.

\noindent 2) IPP avec $u(x)=x$, $u'(x)=1$, $v'(x)=\mathrm{e}^{\,1-x}$, $v(x)=-\mathrm{e}^{\,1-x}$ :
\[\int_0^{\lambda}x\,\mathrm{e}^{\,1-x}\dd x=\big[-x\,\mathrm{e}^{\,1-x}\big]_0^{\lambda}+\int_0^{\lambda}\mathrm{e}^{\,1-x}\dd x
=-\lambda\,\mathrm{e}^{\,1-\lambda}+\big[-\mathrm{e}^{\,1-x}\big]_0^{\lambda}
=-\lambda\,\mathrm{e}^{\,1-\lambda}-\mathrm{e}^{\,1-\lambda}+\mathrm{e}.\]
D'où $\displaystyle\int_0^{\lambda}f(x)\dd x=\mathrm{e}-(\lambda+1)\,\mathrm{e}^{\,1-\lambda}$. Comme
$\lim\limits_{\lambda\to+\infty}(\lambda+1)\,\mathrm{e}^{\,1-\lambda}=\mathrm{e}\lim\limits_{\lambda\to+\infty}(\lambda+1)\mathrm{e}^{-\lambda}=0$,
on obtient $\displaystyle\lim_{\lambda\to+\infty}\int_0^{\lambda}f(x)\dd x=\mathrm{e}$.

\noindent 3) Pour $x\in[0\,;\,1]$ et tout $n\ge0$ : $x^{\,n+1}\le x^{\,n}$ (car $0\le x\le1$), donc
$x^{\,n+1}\mathrm{e}^{-x}\le x^{\,n}\mathrm{e}^{-x}$. Par croissance de l'intégrale sur $[0\,;\,1]$,
$I_{n+1}\le I_n$ : la suite $(I_n)$ est décroissante.

\noindent 4) Sur $[0\,;\,1]$, $0\le\mathrm{e}^{-x}\le1$, donc $0\le x^{\,n}\mathrm{e}^{-x}\le x^{\,n}$. En
intégrant entre $0$ et $1$ : $0\le I_n\le\displaystyle\int_0^1 x^{\,n}\dd x=\Big[\dfrac{x^{\,n+1}}{n+1}\Big]_0^1=\dfrac{1}{n+1}$.
Comme $\dfrac{1}{n+1}\to0$, le théorème des gendarmes donne $\lim\limits_{n\to+\infty}I_n=0$.

\medskip
\noindent\textbf{Exercice 4.}
1) $\Delta=(2\sqrt3)^2-4\times4=12-16=-4=(2i)^2$. Les racines sont
$z=\dfrac{2\sqrt3\pm2i}{2}=\sqrt3\pm i$. Donc $z_A=\sqrt3+i$ et $z_B=\sqrt3-i$.

\noindent 2) $|z_A|=\sqrt{3+1}=2$ et $\arg(z_A)=\dfrac\pi6$ (car $\cos=\frac{\sqrt3}{2}$, $\sin=\frac12$),
donc $z_A=2\big(\cos\frac\pi6+i\sin\frac\pi6\big)=2\mathrm{e}^{\,i\pi/6}$ et $z_B=\overline{z_A}=2\mathrm{e}^{-i\pi/6}$.
Alors $z_A^{\,6}=2^6\,\mathrm{e}^{\,i\pi}=64\times(-1)=-64$.

\noindent 3) $OA=|z_A|=2$, $OB=|z_B|=2$ et $AB=|z_A-z_B|=|2i|=2$. Les trois côtés sont égaux à $2$ :
le triangle $OAB$ est équilatéral.

\noindent 4) $S(z)=2\,\mathrm{e}^{\,i\pi/3}\,z=2\big(\tfrac12+i\tfrac{\sqrt3}{2}\big)z=(1+i\sqrt3)\,z$.
Image de $A$ : $(1+i\sqrt3)(\sqrt3+i)=\sqrt3+i+3i+i^2\sqrt3=(\sqrt3-\sqrt3)+i(1+3)=4i$.
L'affixe de l'image de $A$ est $\mathbf{4i}$.

\noindent 5) Le plan $(MNQ)$ a pour équation par les intercepts $\dfrac x2+\dfrac y3+\dfrac z6=1$,
soit, en multipliant par $6$ : $3x+2y+z=6$. Un vecteur normal est $\vec n(3\,;\,2\,;\,1)$, $\|\vec n\|=\sqrt{14}$.
La distance de $O$ au plan vaut $d(O,(MNQ))=\dfrac{|3\cdot0+2\cdot0+0-6|}{\sqrt{14}}=\dfrac{6}{\sqrt{14}}=\dfrac{3\sqrt{14}}{7}\approx1{,}60\ \text{m}$.

\smallskip
\noindent\textbf{Partie B — Situation.}
\emph{Tâche 1.} La contrainte $2h+b=12$ donne $b=12-2h$ avec $0<h<6$. L'aire de la section est
$A(h)=b\,h=h(12-2h)=12h-2h^2$. Alors $A'(h)=12-4h$, qui s'annule en $h=3$ en changeant de signe
($+$ puis $-$) : $A$ est maximale en $h=3$. Dimensions : $h=\SI{3}{\metre}$ et $b=12-6=\SI{6}{\metre}$ ;
aire maximale $A(3)=3\times6=\SI{18}{\metre\squared}$.

\emph{Tâche 2.} Le volume l'année $n$ suit une suite géométrique de premier terme $V_1=3{,}6$ et de
raison $1{,}12$ : $V_n=3{,}6\times1{,}12^{\,n-1}$. On résout $V_n\ge6$, soit $1{,}12^{\,n-1}\ge\dfrac{6}{3{,}6}=\dfrac53$,
d'où $n-1\ge\dfrac{\ln(5/3)}{\ln1{,}12}\approx\dfrac{0{,}511}{0{,}113}\approx4{,}51$, donc $n\ge6$ (car $n$ entier).
Vérification : $V_5\approx\SI{5,67}{\giga\metre\cubed}<6$ et $V_6\approx\SI{6,34}{\giga\metre\cubed}\ge6$.
La cote d'exploitation est atteinte à partir de la \textbf{6\textsuperscript{e} année}.

\emph{Tâche 3.} Le nombre $T$ de turbines défaillantes suit la loi binomiale $\mathcal{B}(5\,;\,0{,}1)$.
$P(T\ge1)=1-P(T=0)=1-(0{,}9)^5\approx1-0{,}590=\mathbf{0{,}410}$, soit environ \SI{41}{\percent}.
Le nombre moyen attendu est $E(T)=np=5\times0{,}1=0{,}5$ turbine défaillante par an.
